\section{Simultaneous control of all paired owners above the block length}
\label{apx-section}

We now remove the choice of one owner per prime for a precisely specified
paired part of the original boundary. This does not supply a partner for
an isolated high-depth occurrence. Retain the original prime \(n>324\),
\(B=n^4\), \(d=3n-1\), and the integer boundary \(T_k=3F_n(k)\).
For a nonzero integer \(h\) with \(|h|<B\), write
\[
 I_h=\{k\in\mathbb Z:B\le k<2B,\ B\le k+h<2B\},\qquad
 b_q(k,h)=\min\{v_q(T_k),v_q(T_{k+h})\}.
\]
Put \(b_q(h)=\max_{k\in I_h}b_q(k,h)\) and
\begin{equation}
 C_T(h)=\mathop{\mathrm{lcm}}_{k\in I_h}\gcd(T_k,T_{k+h}),\qquad
 C_n(h)=2dn\log64+d^2\log(1+|h|)+\log3.
 \label{apx-constants}
\end{equation}
The shifted-resultant theorem in the preceding section proves
\(\log C_T(h)\le C_n(h)\). The interval \(I_h\) is nonempty,
and every prime sum below has finite support on actual integer boundaries.

\begin{theorem}[The complete sum over paired owners]
\label{apx-all-owners}
Define the actual nonnegative quantity
\[
 L_h^{>}(k)=\sum_{\substack{q>B\\q\ \mathrm{prime}}}
                 \max\{b_q(k,h)-3,0\}\log q .
\]
Then
\begin{equation}
 \sum_{k\in I_h}L_h^{>}(k)
 \le d\sum_{q>B}\max\{b_q(h)-3,0\}\log q
 \le d\log C_T(h)\le dC_n(h).
 \label{apx-total}
\end{equation}
\end{theorem}
\begin{proof}
The leading coefficient of \(F_n\) is \(n3^{3n-2}\).
A prime \(q>B\) divides neither \(n\) nor three, so the reduction
of \(F_n\) modulo \(q\) retains degree \(d\) and has at most
\(d\) roots in \(\mathbb F_q\). The \(B\) consecutive integers
in \([B,2B)\) have distinct residues modulo \(q\). Consequently
at most \(d\) original indices have \(q\mid T_k\).
Imposing the additional hit at \(k+h\) only decreases this number.
For each \(q\), at most \(d\) summands are nonzero and each is
at most \(\max\{b_q(h)-3,0\}\log q\). Interchanging the finite
sums proves the first inequality. The exact exponent of \(q\) in
the lcm \(C_T(h)\) is \(b_q(h)\). Termwise comparison proves
the second inequality; the shifted-resultant bound proves the last.
\end{proof}

\begin{corollary}[One exceptional set for all these primes]
\label{apx-exceptional}
For \(t_k=n\log|3k+\zeta|\) and \(\varepsilon>0\), define
\[
 E_h(\varepsilon)=\{k\in I_h:L_h^{>}(k)>\varepsilon t_k\}.
\]
Then
\begin{equation}
 |E_h(\varepsilon)|<
       \frac{dC_n(h)}{\varepsilon n\log(3B)}.
 \label{apx-markov}
\end{equation}
In particular, for fixed nonzero \(h\),
\[
 |E_h(1/n)|=O_h(n^3/\log n)=o(B).
\]
At every index outside this set the paired quantity is at most
\(t_k/n\). This exceptional set does not depend on choosing a
prime-to-owner assignment.
\end{corollary}
\begin{proof}
Throughout the block,
\(|3k+\zeta|^2=9k^2+3k+1>(3B)^2\), so
\(t_k\ge n\log(3B)\). Sum the strict inequalities defining the
exceptional set and use~\eqref{apx-total}. If the set is empty,
the strict bound still holds because the right side is positive.
For fixed \(h\), \(C_n(h)=O_h(n^2)\) and \(d=O(n)\);
substituting \(\varepsilon=1/n\) gives the asserted estimate.
\end{proof}

\begin{corollary}[Finitely many prescribed shifts]
\label{apx-finite-shifts}
Let \(H\) be a finite set of nonzero integer shifts with \(|h|<B\).
At each original block index take the maximum of \(L_h^{>}(k)\)
over its valid shifts, using zero when none is valid. The number of
indices where this maximum exceeds \(\varepsilon t_k\) is at most
\[
 \frac{d\sum_{h\in H}C_n(h)}{\varepsilon n\log(3B)}.
\]
The bound is strict when \(H\ne\varnothing\). If \(H=\varnothing\),
the maximum and the exceptional set are both zero.
\end{corollary}
\begin{proof}
The maximum is bounded by the sum over valid shifts. Summing
\eqref{apx-total} and applying the preceding counting argument gives
the result. Positivity of the displayed numerator when \(H\ne
\varnothing\) handles the empty exceptional set. No strict inequality
is asserted for an empty shift set.
\end{proof}

\paragraph{Exact scope and the remaining signed contribution.}
For one fixed shift, exactly \(|h|\) endpoints lack a valid partner
inside the block; they must be accounted for in a whole-block assertion.
More seriously, no lower bound has been proved on the fraction of the
full positive prime-depth contribution captured by \(L_h^{>}\).
A prime may have large \(v_q(T_k)\) and no other hit at the required
depth. All excess above the partner's depth remains outside this sum.
The earlier very-large-prime independent-domain theorem even bounds
the common depth of two such independent endpoints by three. Allowing
other endpoints here establishes neither their existence nor their depth.

The theorem bounds a specified positive paired portion. The full signed
identity from the preceding section, including negative depth-one and
depth-two contributions on their original supports, remains unchanged.
The factor \(d\) pays for counting a prime at every actual owner.
For \(q\le B\), residues can repeat and the stated \(d\)-owner argument
does not apply. There is no unproved interchange of a growing shift set
with a fixed-shift estimate, no coverage of isolated primes, and no
exceptional-root or arbitrary-ABC conclusion in these statements.
